\section{Analytic relationships}
We now turn to deriving the exact relationship between:
the distribution of \emph{total} durations selling sex
among the source population: $x \sim p(x)$, and
the distribution of \emph{observed} durations selling sex
among participants in a cross-sectional survey: $z \sim p(z|s)$,
given survey participation status $s$.
We also consider two survey sampling cases:
an ideal case which is representative of the active population, and
a more realistic case which undersamples individuals in the early years of sex work. %SM: suggest "more plausible" instead of "more realistic"?
For reference throughout this section,
Figure~\ref{fig:toy.distr.1} illustrates
example distributions for $p(x)$, $p(z|s)$, and intermediates thereof.
\begin{figure}
  \includegraphics[width=\linewidth]{toy.distr.1}
  \caption{Illustrative distributions of total and censored durations
    among the source, active, and sampled populations
    for 3 coefficients of variation $\CV[x]$ (rows)}
  \label{fig:toy.distr.1}
  \floatfoot{Assumed: $\E[x] = 5$ and $x \sim \txn{Weibull}$
    (Figure~\ref{fig:toy.distr} illustrates other distribution families);
    probability of being sampled per \eqref{eq:p(s|z)}.}
\end{figure}
% ------------------------------------------------------------------------------
\paragraph{Active population bias}
occurs because individuals who sell sex for longer %SM: could tigthen up text by defining/explaining the cause of the bias with the first sentence? [also, if others have used this term or a related concept for other things, suggest citing here?]
are more likely to be active at any given time.
Specifically, the probability of an individual being active
is proportional to their total duration $x$.
Thus, the distribution of total durations among the active population is:
\begin{equation}
  p(x|a) = \norm ~ x ~ p(x)
  \label{eq:p(x|a)}
  \enlargethispage{1\baselineskip} % TEMP
\end{equation}
where $K$ is a normalizing constant.
These total durations among the active population are typically unobserved,
except in cohort studies which follow participants until they exit sex work.
From \eqref{eq:p(x|a)} we can see that
the mean total duration among the active population $\E[x|a]$
is always larger than the mean among the source population $\E[x]$
(Figure~\ref{fig:toy.distr}, light green vs gray circles);
specifically, we have that:
\begin{equation}
  \E[x|a] = \E[x] \left(1 + \CV^2[x]\right)
  \label{eq:E[x|a]}
\end{equation}
This is because greater variability in durations among the source population,
as measured by the coefficient of variation: $\CV[x]$,
allows those with longer durations to increasingly dominate the active population.
We also note that the probability of observing
active individuals with very short total durations is inherently very small
(Figure~\ref{fig:toy.distr.1}, light green).
% ------------------------------------------------------------------------------
\paragraph{Censoring}
occurs because with random timing of the survey,
observed durations at the time of survey ($z$) will be right-censored %SM: minor edit for considration - same rationale as above
and uniformly distributed between 0 and the total duration ($x$):
\begin{equation}
  p(z|x) = \txn{Unif}(0,x) =
  \begin{cases}
    \frac{1}{x} & \txn{if}~ 0 \le z \le x \\
    0 & \txn{else}
  \end{cases}
  \label{eq:p(z|x)}
\end{equation}
Since the mean of $\txn{Unif}(0,x)$ is $\frac{x}{2}$,
the mean censored duration is always half the mean total duration
among the active population: $\E[z|a] = \frac{1}{2} \E[x|a]$
(Figure~\ref{fig:toy.distr.1}, dark green vs light green circles).
This idea appears to be well-known, and likely motivates  %SM: clarify what is meant here? 
doubling the mean $\E[z|a]$ to estimate $\E[x]$ in method~3 of \cite{Fazito2012}.
However, method~3 does not account for selection of the active population noted above,
and thus actually estimates the mean total duration among the active population $\E[x|a]$,
which we know to be greater than among the source population $\E[x]$ per \eqref{eq:E[x|a]}.%
\footnote{\citet{Fazito2012} partially acknowledge this issue
  by opting to use the median over the mean of $z|a$ where possible,
  since the median is less sensitive to outliers.
  However, as shown in Figures \ref{fig:toy.distr} and \ref{fig:toy.est},
  doubling the median $\M[z|a]$ still overestimates $\E[x]$,
  just to a lesser degree than doubling the mean $\E[z|a]$,
  since the median will always be less than the mean for right-skewed distributions.}
\par
If total durations among the source population are exponentially distributed: %SM: nice
$p(x) \sim \txn{Exp}(\lambda)$, then $\CV[x] = 1$, and so
the mean \emph{censored} duration among the \emph{active} population will directly equal
the mean \emph{total} duration among the \emph{source} population (without doubling):
$\E[z|a] = \E[x]$.
In other words, under these conditions,
the mean reported duration in a typical cross-sectional survey
will be an unbiased estimator of the target $\E[x]$.
Moreover, even the \emph{distributions} will be identical:
$p(z|a) = p(x) = \txn{Exp}(\lambda)$, as we will show below.
This is particularly notable since
compartmental models usually assume exponentially distributed durations.
% ------------------------------------------------------------------------------
\paragraph{Complete distribution}
By combining \eqrefa{eq:p(x|a)}{eq:p(z|x)}, we can derive
the complete distribution of censored durations among the active population $p(z|a)$
from the distribution of total durations among the source population $p(x)$
\cite{Tufto2017}:
\begin{equation}
  \begin{aligned}
    p(z|a) &= \int_0^\infty p(z|x) ~ p(x|a) ~ dx \\
           &= \int_z^\infty \frac{1}{x} ~ p(x|a) ~ dx \\
           &= \norm ~ \int_z^\infty \frac{x}{x} ~ p(x) ~ dx \\
           &= \norm ~ \left(1 - F_x(z)\right)
  \end{aligned}
  \label{eq:p(z|a)}
  \enlargethispage{1\baselineskip} % TEMP
\end{equation}
Interestingly, this is simply %SM: style-edit: consider removing "interestingly", and instead something like. It turns out that this is simply one minus the cumulative.... [but am not fussed - see if similar language/framing in other practice of epi papers, just to make it eadier for editor to see that the style matches something they are used to.
one minus the cumulative distribution of total durations $F_x$,
normalized to have unit area under the curve (via $K$).%
\footnote{The integral lower bound changes from 0 to $z$ because
  $p(z|x) = \frac{1}{x}$ only when $x < z$ and $p(z|x) = 0$ otherwise,
  per \eqref{eq:p(z|x)}.}
For the exponential case, we have:
$F_x = 1 - e^{-\lambda x}$, $K = \lambda^{-1}$, and thus
$p(z|a) = \lambda e^{-\lambda x} = p(x)$.
\par
More generally, $\CV[z|a]$ is typically drawn towards 1 relative to $\CV[x]$,
attenuating skewness when $\CV[x] > 1$, and
amplifying skewness when $\CV[x] < 1$
(Figures \ref{fig:toy.distr.1} and \ref{fig:toy.distr}).
However, we believe that
$\CV[x] > 1$ is more likely than $\CV[x] < 1$ for two reasons.
First, heavy-tailed distributions having $\CV[x] > 1$
arise naturally via selection effects
in populations with heterogeneous rates \cite{Balan2020}.
Second, most parametric distributions with $\CV[x] < 1$
have $p(x\sp\rarr0)\sp\rarr0$ (Figure~\ref{fig:toy.distr}),
which effectively implies
that few people exit sex work shortly after starting, or
that most people stay in sex work for a "minimum" period,
which seems unrealistic.
If $\CV[x] > 1$ is indeed more likely, then
the skewness of the observable duration distribution $p(z|a)$
will generally underestimate
the skewness of the true duration distribution $p(x)$.
% ------------------------------------------------------------------------------
\paragraph{Sampling Biases}
If the survey undersamples individuals in the early years of sex work,
this probability of being sampled will be a function of
current duration ($z$) not total duration ($x$),
so we denote it $p(s|z)$.
Many parameterizations are possible; \eg \eqref{eq:p(s|z)}.
The distribution of observed durations among survey participants
is then updated from \eqref{eq:p(z|a)} simply as:
\begin{equation}
  p(z|s) = \norm ~ p(s|z) \left(1 - F_x(z)\right)
  \label{eq:p(z|s)}
\end{equation}
% ------------------------------------------------------------------------------
\paragraph{Estimating $\E[x]$}  %SM: this is really nice
We conclude this section by revisiting
the three methods considered by \citet{Fazito2012}
to estimate the "true" mean duration in sex work $\E[x]$.
Figure~\ref{fig:toy.est} plots:
the typically unobserved mean and median of total durations (method~1),
\emph{double} the means and medians of censored durations (method~3),
and the reciprocal proportion of censored durations less than 1 year (method~2),
across a range of $\CV[x]$ and distribution families.
% JK: focusing on CV[x] makes the trends more straightforward,
%     but CV[x] is unknown in reality.  %SM: this is worth "testing" out with a reader - I like and understand it, as is, with explanation in the discussion section. but could be worth testing it out with Linwei or (or maybe even Caroline?) as a "reader/peer-reviewer" for this part? 
%     I wonder if we should consider changing the x-axis to
%     CV[*] where * is the data used to calculate the estimate?
%     One challenge is that for many distributions, CV[z|a] > ~ 1.2
%     is impossible to achieve when truncating at 50 years max duration.
%     I was experimenting with some proxies for CV[x] based on p(z|a),
%     but they are all rather sensitive to the truncation point (50)
%     and/or assumed distribution family (Gamma, etc.).
We see that at low $\CV[x]$, nearly all methods converge on $\E[x]$.
However, as $\CV[x]$ increases,
selection towards longer total durations among the active/sampled populations %SM: to clarify, active or sampled, or both? 
raises the means of the corresponding censored durations $\E[z|\cdot]$,
according to \eqref{eq:E[x|a]}.
Median censored durations exhibit a similar trend,
but lag below means (true for all right-skewed distributions)
to a degree which depends on the family,
making such medians less useful estimators of $\E[x]$.
For all method~3 estimates, undersampling of early sex work
(here assuming $\pp = 0.5, \tt = 1$ year)
predictably biases estimates slightly higher;
this bias is also larger for smaller $\E[x]$ (not shown).
Finally, turnover-based estimates (method~2) tends to be most consistent,
only slightly overestimating $\E[x]$ as $\CV[x]$ increases;
however, this method is especially sensitive to
undersampling of early sex work, which biases estimates upwards.
\par
Nevertheless, the best simple estimate of $\E[x]$ is likely $\E[z|s]$,
assuming $\CV[x] \approx 1$, or more generally:
\begin{equation}
  \sE[x] = \frac{1}{2} ~ \E[z|s] \left(1 + \CV^2[x]\right)
\end{equation}
